\documentclass{article}
\usepackage[paperwidth=240mm,paperheight=115mm,margin=12mm]{geometry}
\usepackage{unicode-math,xcolor,array}
\setmathfont{XITS Math}
\pagestyle{empty}\setlength\parindent{0pt}
\newbox\squarebody
\newdimen\axis\newdimen\radius
\newcommand\LeftSized{\mathopen{\nulldelimiterspace=0pt\left[\vphantom{\vrule width0pt height\dimexpr\radius+\axis\relax depth\dimexpr\radius-\axis\relax}\right.}}
\newcommand\RightSized{\mathclose{\nulldelimiterspace=0pt\left.\vphantom{\vrule width0pt height\dimexpr\radius+\axis\relax depth\dimexpr\radius-\axis\relax}\right]}}
\newcommand\Trial[1]{\directlua{
local r=math.max(tex.box.squarebody.height-tex.dimen.axis,tex.box.squarebody.depth+tex.dimen.axis)/65536
local W=tex.box.squarebody.width/65536
local lambda=W/2
local integral=r^(#1)*lambda*(1-math.exp(-W/lambda))
local norm=lambda*(1-math.exp(-W/lambda))
local result=(integral/norm)^(1/(#1))
assert(math.abs(result-r)<1e-10)
tex.dimen.radius=math.floor(result*65536+.5)
}\centering\strut $q=#1$\par\bigskip
$\displaystyle\LeftSized\copy\squarebody\RightSized$\par\bigskip
{\small $R=10.00$ pt}\par}
\begin{document}
\axis=\Umathaxis\displaystyle
\setbox\squarebody=\hbox{\vrule width20pt height\dimexpr10pt+\axis\relax depth\dimexpr10pt-\axis\relax}
{\Large\bfseries One square: native delimiters versus integral means}\par\medskip
A single solid square, $20\,\mathrm{pt}\times20\,\mathrm{pt}$, centered on the math axis.
No gaps or other content.\par\bigskip
\begin{tabular}{@{}p{36mm}p{36mm}p{36mm}p{36mm}p{36mm}@{}}
\centering\strut Native\par\bigskip
$\displaystyle\left[\copy\squarebody\right]$\par\bigskip {\small\texttt{\string\left...\string\right}}\par
&\Trial{1}&\Trial{2}&\Trial{8}&\Trial{16}
\end{tabular}\par\bigskip
$\lambda=W/2$, $p=3$. Both side estimates equal $10$ pt for every $q$.\par\smallskip
The same native font and delimiter parameters are used throughout.
\end{document}
